
% verifique abaixo e em tutoriais latex o uso de itemize, enumerate, pacotes e ambientes de tabelas e figuras, citação (normal e apud), etc


%\begin{comment}
\chapter{Utilidades}\label{cap:utilidades}
Esse capítulo mostrará como inserir itens, figuras, equações e tabelas, com exemplos. Além disso, mostrará como citar referências.

\section{Inserindo...}

\subsection {Itens}
Para inserir itens, dentro das tags $\backslash$begin\{itemize\} e $\backslash$end\{itemize\}, colocá-los dentro de $\backslash$item \{\}, como no exemplo abaixo. O parâmetro $\backslash$itemsep indica o espaçamento entre os itens e o parâmetro $\backslash$itemindent indica a indentação do texto dos itens.



\begin{verbatim}
    \begin{itemize}
    \setlength{\itemsep}{1pt}
    \setlength{\itemindent}{20pt}
         \item{ item 1}
         \item{ item 2}
         \item{ item 3}
    \end{itemize}
\end{verbatim}

O resultado é o seguinte:
\begin{itemize}
\setlength{\itemsep}{1pt}
\setlength{\itemindent}{20pt}
     \item{ item 1}
     \item{ item 2}
     \item{ item 3}
\end{itemize}

\subsection {Figuras}
Para inserir uma figura, dizer o nome e a localização da mesma dentro do \{\} no final da linha  $\backslash$includegraphics. O parâmetro [ht] indica que a posição provável que a figura estará será no lugar que ela foi inserida (here) e no topo. É possível mudar para [hb], por exemplo, se desejar que a figura fique o mais próximo da posição inserida, mas no final da página (bottom). O resultado é a Figura \ref{fig:ic}.
\begin{verbatim}
\begin{figure*}[ht]
 \centering
 \includegraphics[width=0.30\textwidth,natwidth=227,natheight=200]{ic-uff.jpg}
 \caption{\label{fig:ic}Instituto de Computação}
\end{figure*}
\end{verbatim}



\subsection{Equações}
É possível escrever diversos tipos de equações. Para inserir uma equação, é preciso colocá-la dentro das tags $\backslash$begin\{equation\} e $\backslash$end\{equation\}. Exemplo de código a seguir:

\begin{verbatim}
    \begin {equation}
           \overline{d}=\frac{1}{n} \sum\limits_{i\in V}d_i
    \end {equation}
\end{verbatim}

O código acima resulta na seguinte equação com fração e somatório:
\begin {equation}
\overline{d}=\frac{1}{n} \sum\limits_{i\in V}d_i
\end {equation}

\subsection {Tabelas}
O código a seguir mostra como inserir uma tabela. Cada `c'' depois de $\backslash$begin\{tabular\} corresponde a uma coluna. Assim, se a tabela tiver 4 colunas, no código devem ter 4 c's. As colunas são separadas por \&. O resultado é a Tabela \ref{tab:tabela1}.

\begin{verbatim}
\begin{table}[ht]
\centering
\scriptsize
\begin{tabular}{cccc}
\hline
$i-j$ & $g_{i4j}$ & $g_{ij}$ & $\frac{g_{i4j}}{g_{ij}}$\\
\hline
1-2 & 1 & 2 & 0.5 \\
2-3 & 1 & 2 & 0.5 \\
3-6 & 3 & 3 & 1 \\
\hline
Betweenness do& vértice 4& = &2\\
\hline
\end{tabular}
        \caption{\label{tab:tabela1} Legenda da Tabela}
\end{table}
\end{verbatim}

\begin{table}[ht]
\centering
\begin{tabular}{cccc}
\hline
$i-j$ & $g_{i4j}$ & $g_{ij}$ & $\frac{g_{i4j}}{g_{ij}}$\\
\hline
1-2 & 1 & 2 & 0.5 \\
1-3 & 0 & 1 & 0 \\
1-6 & 1 & 1 & 1 \\
2-3 & 1 & 2 & 0.5 \\
5-6 & 3 & 3 & 1 \\
\hline
$Betweenness$ do& vértice 4& = &3\\
\hline
\end{tabular}
        \caption{\label{tab:tabela1} Legenda da Tabela}
\end{table}



\section{Arquivo da bibliografia}

Para usar o recurso do BibTeX no Latex é preciso editar um arquivo
bibliografia.bib, com as informações sobre a bibliografia.
Apresenta-se a seguir um exemplo de para referenciar um livro:

\bigskip

@book\{kant,

  title=\{Critica da razão pura\},

  author=\{Kant, I.\},

  year=\{1781\},

  publisher=\{Abril Cultural São Paulo\}


Uma referência genérica pode ser escrita neste arquivo como:

\bigskip


@TIPO\{ApelidoDaReferencia,

AUTHOR = "Fulano de Tal",

TITLE = "O Título da  trabalho",

a PUBLISHER = " Editora",

ADDRESS = "Cidade, etc",

. . . YEAR = "1998" \}

\bigskip
\noindent onde o TIPO pode ser, por exemplo, BOOK (para livros),
ARTICLE (para artigos), INCOLLECTION (para partes de um livro),
PROCEEDINGS (para artigo em atas de congressos) ou TECHREPORT (para
relatórios técnicos).  Cada TIPO tem um conjunto de campos
obrigatório e outro conjunto de campos que é opcional. Na dúvida,
consulte um tutorial como já sugerido neste documento.

\subsection{Citando referências}

As referências devem ser citadas ao longo do texto pelo número. Alguns exemplos estão a seguir.

Esse template está disponível no site do Curso de Ciência da Computação \cite{IC}''. Se desejar citar um artigo publicado em periódico, cite conforme definito em \cite{Rocha:NomeDoPeriodico2013}. Já um artigo publicado em conferência, referencie como feito em \cite{Franca:NomeDaConferencia2012}. Agora, se a citação for de um livro, use \cite{Toledo:NomeDoLivro2011}. Em alguns momentos, será apropriado introduzir ainda o nome dos autores no texto, como poderia ter sido feito por Toledo el al em \cite{Toledo:NomeDoLivro2011} ou se fossemos referenciar a dissertação de mestrado da Cinthya França \cite{Franca:MSc2013}.

Para cada referência a ser citada, se faz necessário acrescentá-la no arquivo biblioteca.bib. É preciso compilar o arquivo monografiamodelo.tex mais de uma vez, pois na primeira o compilador lê as referências e cria os identificadores e, na segunda, ele atribui a numeração correta.

Para quem está vendo LaTex pela primeira vez, é indicado dar uma olhada em alguns tutoriais. A definição do Wikipedia aponta para uma série de tutoriais disponíveis na Web \cite{WikipediaLatex}.
%\end{comment}